\begin{abstract}
This study evaluates the performance of three domain-specific computer-vision tools (AIRVIC, DVICE, and Cellpose) and four general-purpose multimodal large language models (ChatGPT, Claude, Gemini, and Grok) on the task of detecting and classifying cytopathic effect (CPE) in Vero-cell culture microscope images. Ground-truth CPE presence and morphological types were derived solely from narrative descriptions supplied by the contract research organization (CRO) that performed the cell-culture experiments and were cross-referenced against the standardized CPE descriptors listed in Table 7 of the CLSI M41-A Guideline. Quantitative analysis was performed exclusively on the 22 images that received CRO descriptions.

None of the evaluated tools demonstrated adequate performance for reliable CPE detection. Domain-specific tools showed extremely high false-positive rates, while the multimodal models exhibited either strong conservative bias or only modest balanced accuracy. These findings indicate that current AI approaches, whether domain-specific or general-purpose, are not suitable as objective tools for automated CPE analysis of the Vero cell culture image dataset used in this study.
\end{abstract}